\documentclass{beamer}

\usepackage[italian]{babel}
\usepackage[utf8]{inputenc}
\usepackage{listings}
\usepackage{xcolor}

\usetheme{Madrid}

\title{Test di Web Service REST}
\subtitle{Uso di Curl e Postman}
\author{Prof. Fedeli Massimo }
\date{}

\definecolor{codegray}{rgb}{0.95,0.95,0.95}

\lstset{
	backgroundcolor=\color{codegray},
	basicstyle=\ttfamily\small,
	breaklines=true
}

\begin{document}
	
	\frame{\titlepage}
	
	%------------------------------------------------
	
	\begin{frame}{Web Service REST}
		Un \textbf{Web Service REST} è un servizio web che utilizza il protocollo HTTP per permettere la comunicazione tra applicazioni.
		
		\bigskip
		
		Caratteristiche principali:
		
		\begin{itemize}
			\item Architettura client-server
			\item Comunicazione tramite HTTP
			\item Utilizzo di URI per identificare le risorse
			\item Dati spesso rappresentati in formato JSON
		\end{itemize}
		
	\end{frame}
	
	%------------------------------------------------
	
	\begin{frame}{Operazioni REST}
		Le operazioni sui dati vengono effettuate tramite i metodi HTTP.
		
		\bigskip
		
		\begin{center}
			\begin{tabular}{c c c}
				\textbf{Metodo} & \textbf{Operazione} & \textbf{Descrizione} \\
				\hline
				GET & Read & Recupera dati \\
				POST & Create & Inserisce nuovi dati \\
				PUT & Update & Aggiorna dati \\
				DELETE & Delete & Cancella dati
			\end{tabular}
		\end{center}
		
	\end{frame}
	
	%------------------------------------------------
	
	\begin{frame}{Strumenti per testare API REST}
		Durante lo sviluppo di un web service è necessario verificare il corretto funzionamento delle API.
		
		\bigskip
		
		Strumenti comuni:
		
		\begin{itemize}
			\item Curl (strumento a riga di comando)
			\item Postman (interfaccia grafica)
		\end{itemize}
		
		\bigskip
		
		Questi strumenti permettono di simulare richieste HTTP verso il server.
		
	\end{frame}
	
	%------------------------------------------------
	
	\begin{frame}{Il comando Curl}
		\textbf{Curl} è un software open source utilizzato per trasferire dati tramite URL.
		
		\bigskip
		
		Caratteristiche:
		
		\begin{itemize}
			\item disponibile su Windows, Linux e MacOS
			\item utilizzabile da terminale
			\item permette di inviare richieste HTTP
		\end{itemize}
		
	\end{frame}
	
	%------------------------------------------------
	
	\begin{frame}[fragile]{Verifica installazione Curl}
		
		Nel terminale digitiamo:
		
		\begin{lstlisting}
			curl --version
		\end{lstlisting}
		
		Il comando mostra la versione installata e i protocolli supportati.
		
	\end{frame}
	
	%------------------------------------------------
	
	\begin{frame}{Web Service di esempio}
		Supponiamo di avere il servizio:
		
		\begin{center}
			http://localhost:8080/PizzaExpress/api/pizze
		\end{center}
		
		Il servizio gestisce un database di pizze con i seguenti campi:
		
		\begin{itemize}
			\item id
			\item nome
			\item prezzo
		\end{itemize}
		
	\end{frame}
	
	%------------------------------------------------
	
	\begin{frame}[fragile]{Metodo GET}
		
		Il metodo \textbf{GET} serve per recuperare le risorse.
		
		\bigskip
		
		\begin{lstlisting}
			curl http://localhost:8080/PizzaExpress/api/pizze
		\end{lstlisting}
		
	\end{frame}
	
	%------------------------------------------------
	
	\begin{frame}[fragile]{Risposta del server}
		
		Il server restituisce dati in formato JSON.
		
		\begin{lstlisting}
			{
				"pizze": [
				{"id":1, "nome":"Margherita", "prezzo":4.5},
				{"id":2, "nome":"Prosciutto", "prezzo":6.5},
				{"id":3, "nome":"Capricciosa", "prezzo":8.0}
				]
			}
		\end{lstlisting}
		
	\end{frame}
	
	%------------------------------------------------
	
	\begin{frame}{Messaggi HTTP}
		
		Ogni risposta HTTP contiene uno status code.
		
		\bigskip
		
		Codici comuni:
		
		\begin{itemize}
			\item 200 OK
			\item 201 Created
			\item 204 No Content
			\item 404 Not Found
		\end{itemize}
		
	\end{frame}
	
	%------------------------------------------------
	
	\begin{frame}[fragile]{Metodo PUT}
		
		Il metodo \textbf{PUT} permette di aggiornare una risorsa.
		
		Esempio: aggiornare il prezzo della pizza con id 1.
		
		\begin{lstlisting}
			curl -X PUT http://localhost:8080/PizzaExpress/api/pizze/1 \
			-H "Content-Type: application/json" \
			-d '{"prezzo":4.7}'
		\end{lstlisting}
		
	\end{frame}
	
	%------------------------------------------------
	
	\begin{frame}{Risposta del server}
		
		Dopo l'aggiornamento il server restituisce:
		
		\begin{itemize}
			\item codice HTTP 204 (No Content)
		\end{itemize}
		
		Questo indica che la modifica è stata effettuata correttamente.
		
	\end{frame}
	
	%------------------------------------------------
	
	\begin{frame}[fragile]{Metodo POST}
		
		Il metodo \textbf{POST} permette di inserire una nuova risorsa.
		
		\begin{lstlisting}
			curl -X POST http://localhost:8080/PizzaExpress/api/pizze \
			-H "Content-Type: application/json" \
			-d '{"id":4,"nome":"Fantasia","prezzo":8.5}'
		\end{lstlisting}
		
	\end{frame}
	
	%------------------------------------------------
	
	\begin{frame}{Risposta POST}
		
		Se l'inserimento avviene correttamente il server restituisce:
		
		\begin{itemize}
			\item 201 Created
		\end{itemize}
		
		La nuova pizza viene inserita nel database.
		
	\end{frame}
	
	%------------------------------------------------
	
	\begin{frame}[fragile]{Metodo DELETE}
		
		Il metodo DELETE permette di eliminare una risorsa.
		
		\begin{lstlisting}
			curl -X DELETE \
			http://localhost:8080/PizzaExpress/api/pizze/2
		\end{lstlisting}
		
	\end{frame}
	
	%------------------------------------------------
	
	\begin{frame}{Risposta DELETE}
		
		Se la cancellazione ha successo:
		
		\begin{itemize}
			\item HTTP 204 No Content
		\end{itemize}
		
		Se la risorsa non esiste più:
		
		\begin{itemize}
			\item HTTP 404 Not Found
		\end{itemize}
		
	\end{frame}
	
	%------------------------------------------------
	
	\begin{frame}{Introduzione a Postman}
		
		\textbf{Postman} è una piattaforma grafica per testare API.
		
		\bigskip
		
		Permette di:
		
		\begin{itemize}
			\item inviare richieste HTTP
			\item analizzare le risposte
			\item salvare test
			\item organizzare collezioni di richieste
		\end{itemize}
		
	\end{frame}
	
	%------------------------------------------------
	
	\begin{frame}{Richiesta GET con Postman}
		
		Procedura:
		
		\begin{enumerate}
			\item Aprire Postman
			\item Selezionare metodo \textbf{GET}
			\item Inserire l'URL del servizio
			\item Premere Send
		\end{enumerate}
		
	\end{frame}
	
	%------------------------------------------------
	
	\begin{frame}[fragile]{Aggiornamento con PUT}
		
		In Postman:
		
		\begin{itemize}
			\item selezionare metodo \textbf{PUT}
			\item impostare header Content-Type: application/json
			\item inserire dati JSON nel body
		\end{itemize}
		
		\begin{lstlisting}
			{
				"prezzo": 4.7
			}
		\end{lstlisting}
		
	\end{frame}
	
	%------------------------------------------------
	
	\begin{frame}[fragile]{Inserimento con POST}
		
		Passaggi:
		
		\begin{itemize}
			\item metodo POST
			\item body in formato JSON
		\end{itemize}
		
		\begin{lstlisting}
			{
				"id": 4,
				"nome": "Fantasia",
				"prezzo": 8.5
			}
		\end{lstlisting}
		
	\end{frame}
	
	%------------------------------------------------
	
	\begin{frame}{Eliminazione con DELETE}
		
		Procedura:
		
		\begin{itemize}
			\item metodo DELETE
			\item indicare l'ID della risorsa
		\end{itemize}
		
		\begin{center}
			api/pizze/2
		\end{center}
		
	\end{frame}
	
	%------------------------------------------------
	
	\begin{frame}{Verifica delle modifiche}
		
		Dopo aver eseguito le operazioni:
		
		\begin{itemize}
			\item inviamo nuovamente una richiesta GET
			\item controlliamo i dati restituiti
		\end{itemize}
		
		Questo permette di verificare se le modifiche sono state applicate correttamente.
		
	\end{frame}
	
	%------------------------------------------------
	
	\begin{frame}{Vantaggi di Curl e Postman}
		
		\begin{itemize}
			\item testing rapido delle API
			\item verifica dei codici HTTP
			\item debug dei servizi web
			\item simulazione delle richieste client
		\end{itemize}
		
	\end{frame}
	
	%------------------------------------------------
	
	\begin{frame}{Conclusioni}
		
		Durante lo sviluppo di API REST è fondamentale testare i servizi.
		
		\bigskip
		
		Curl e Postman consentono di:
		
		\begin{itemize}
			\item inviare richieste HTTP
			\item analizzare le risposte
			\item verificare il corretto funzionamento del web service
		\end{itemize}
		
	\end{frame}
	
\end{document}